\documentclass[12pt]{article}
\usepackage[T1]{fontenc}
\usepackage{amsmath,amssymb,amsthm}
\usepackage{hyperref}
\usepackage{geometry}
\usepackage{graphicx}
\geometry{margin=1in}

\newtheorem{definition}{Definition}
\newtheorem{proposition}{Proposition}
\newtheorem{corollary}{Corollary}
\newtheorem{conjecture}{Conjecture}
\newtheorem{remark}{Remark}
\newtheorem{example}{Example}

\title{The Topology of Reliable Cognition:\\
Grounding Structure, Cognitive Distance,\\
and the Dynamics of Abductive Overwrite}

\author{Minamo Minamoto\\
Independent\\
\texttt{minamominamoto4f5683f6@gmail.com}}

\date{April 2026}

\begin{document}

\maketitle

\begin{abstract}
We propose a structural framework linking symbol grounding and abductive
dynamics in cognitive systems.
Rather than treating grounding as a binary property, we model it as
inducing a topology over a shared concept set, and introduce
\emph{cognitive distance} ($d_\mathrm{cog}$) to quantify structural
similarity between agents.

We operationalize $d_\mathrm{cog}$ via representational similarity
analysis (RSA), defining a second-order correlation metric
($\rho_\mathrm{cog}$) over concept-level activation patterns.
Experiments on language models (Pythia-70m training checkpoints and
GPT-2) show that relational structure becomes increasingly aligned
through training, supporting the view that conceptual structure
evolves continuously rather than via abrupt overwrites.
Cross-model replication ($\rho_\mathrm{cog} \approx 0.42$ between
GPT-2 and fully trained Pythia vs.\ $\approx 0.10$ vs.\ random
initialization) suggests that training induces a shared representational
geometry independent of architecture.

We further introduce \emph{abductive overwrite} as a mechanism that
restructures interpretive frames under hypothesis revision, and identify
\emph{negative abductive capability} (NAC) as a structural condition
for maintaining grounding stability.
Evidence from comparative cognition (\textit{C.~elegans},
\textit{Drosophila}, \textit{Apis mellifera}) illustrates that
behavioral equivalence does not imply representational equivalence.
These notions generate three testable predictions linking training
dynamics, representational geometry, and behavioral outcomes.

Our results suggest that grounding should be understood as a structural,
measurable property, providing a bridge between cognitive theory and
empirical analysis of artificial systems.
\end{abstract}

\noindent\textbf{Keywords:} symbol grounding, abduction, cognitive distance, negative abductive capability, representational similarity analysis, comparative cognition

\tableofcontents

\section{Introduction}
% ============================================================

Two questions about cognition have received sustained but largely
separate treatment in the literature.
The first is the \emph{symbol grounding problem}: how do internal
symbols acquire meaning anchored in sensorimotor experience
\cite{harnad1990}?
The second is the \emph{problem of abduction}: how do cognitive
systems generate and revise hypotheses under uncertainty
\cite{peirce1903,harman1965}?

We propose a framework in which these questions are structurally
connected: the topology of the sensory space determines both the
geometry of cognitive representations and the dynamics of hypothesis
revision.

The standard treatment of symbol grounding is binary: either a symbol
system is causally connected to sensorimotor states (grounded) or it
is not.
We argue that this framing conceals the deeper question: grounded
\emph{in what}, and \emph{how}?
Two agents may each satisfy Harnad's condition while inhabiting mutually
incommensurable representational spaces.
The existence of grounding does not imply the equivalence of grounding.

The standard treatment of abduction presents hypothesis revision as
correction: a bad hypothesis is detected and patched while the
underlying representational architecture remains intact.
We argue that this framing is structurally wrong: when an agent adopts
a new hypothesis, that hypothesis restructures the interpretive frame
through which all subsequent input is processed.
Revision is therefore not correction but \emph{overwrite}.
A premature overwrite does not introduce a local error; it establishes
a biased initial condition for all downstream cognition.

Connecting these two claims: the grounding-induced topology of a
cognitive system is not fixed at birth but updated through abductive
transitions.
Premature overwrite is therefore a \emph{dynamic source} of topological
distortion---it increases $d_\mathrm{cog}$ not by mismatching sensory
spaces but by committing to an interpretive frame before the evidence
base is sufficient to determine the correct topology.

Reliable cognition requires both the right structure (\emph{grounding
equivalence}) and the right dynamics (\emph{negative abductive
capability}).
A system that has the right sensory apparatus but commits prematurely
will accumulate topological distortion over time.
A system that suspends commitment appropriately preserves the structural
integrity of its representational space.

\paragraph{Relation to prior work.}
G\"{a}rdenfors's theory of conceptual spaces \cite{gardenfors2000}
provides the closest precedent for a geometric approach to meaning.
Our framework is complementary but distinct: we derive the geometry
from sensory transduction rather than stipulating quality dimensions in
advance, and we use topology rather than Euclidean convexity.
The abduction literature from Peirce onward has framed hypothesis
revision as search within a fixed representational space; our overwrite
model replaces the fixed-space assumption.

\paragraph{Scope and falsifiability.}
This paper is a theoretical proposal, not a completed formal theory.
The central claims are advanced as structural hypotheses supported by
definitional argument and phenomenological evidence from comparative
cognition; their full formal generalization is not established.
We identify three classes of empirical evidence that would bear on the
theory.
First, the $d_\mathrm{cog}$ framework predicts that species sharing
behavioral rules but not sensory topologies will show systematically
higher $d_\mathrm{cog}$ as measured by representational similarity
analysis \cite{kriegeskorte2008}; this is testable via neural recording
in behavioral paradigms that equate transfer accuracy across species.
Second, the NAC hypothesis predicts that agents with higher commitment
thresholds will show lower premature overwrite rates under ambiguous
stimuli with gradually accumulating evidence; this is measurable as
commitment latency and reversibility in both biological and artificial
systems.
Third, the overwrite-as-$d_\mathrm{cog}$-increase hypothesis predicts
that training conditions that suppress suspension (e.g., contrastive
objectives) will produce higher $d_\mathrm{cog}$ relative to an
ideally grounded reference, measurable via RSA on learned
representations.
Open Problems (Section~8) identify the most pressing formalization
and measurement obligations.

\medskip
\noindent
\textbf{Central variables.}
The two central variables of this paper are $d_\mathrm{cog}$ and NAC.
$d_\mathrm{cog}$ is the static measure: how far apart are two agents'
grounding-induced topologies?
NAC is the dynamic condition: does the agent have the capacity to
prevent that gap from widening through premature commitment?
All other constructs in the paper---suspension region, commitment
threshold, overwrite transition, bias propagation factor---are
supporting definitions that make these two variables precise.

% ============================================================
\begin{figure}[t]
\centering
\includegraphics[width=0.82\textwidth]{cogsci_unified_fig1_concept.png}
\caption{
\textbf{Grounding structure and cognitive distance.}
Each agent's sensory space $S(A)$ induces a topology $\tau_A$ on the
shared symbol set $\Sigma$ via the grounding map $G_A$.
The cognitive distance $d_\mathrm{cog}(A,B)$ measures the structural
gap between $\tau_A$ and $\tau_B$, operationalized as the
second-order RSA correlation $\rho_\mathrm{cog}$ (see
Remark~\ref{rem:rsa-operationalization}).
}
\label{fig:concept}
\end{figure}

\section{The Topology of Grounding}
% ============================================================

\subsection{Sensory Spaces}

Let $A$ be a cognitive agent equipped with transduction channels
$\mathcal{T} = \{T_1, \ldots, T_n\}$, where each $T_i$ maps physical
stimuli in domain $\Phi_i$ to internal activation states $A_i$.

\begin{definition}[Sensory Space]
The sensory space of agent $A$ is the product space
\[
S(A) = \prod_{i=1}^{n} A_i
\]
equipped with the product topology, where each $A_i$ carries the
topology induced by the neural activation metric of channel $T_i$.
\end{definition}

\begin{remark}
This definition is intentionally coarse. Cross-modal interactions and
attentional modulation couple the channels; a refined treatment would
replace the product with a fiber bundle. We use the product topology
as a first approximation and flag the refinement as OP2.
\end{remark}

\subsection{Grounding as a Map}

Let $\Sigma$ be a set of symbols. Let $\mathcal{K}(S(A))$ denote
the hyperspace of nonempty closed subsets of $S(A)$, equipped with
the Vietoris topology \cite{vietoris1922}.

\begin{definition}[Weak Grounding]
A grounding map for $A$ is a function $G : \Sigma \to \mathcal{K}(S(A))$
where $G(\sigma)$ is the closed set of sensory states associated with
symbol $\sigma$. The symbol set $\Sigma$ is \emph{weakly grounded}
in $A$ if $G(\sigma) \neq \emptyset$ for all $\sigma \in \Sigma$.
\end{definition}

Weak grounding is Harnad's condition. It rules out purely formal
symbol systems but says nothing about the structure of $G(\sigma)$
within $S(A)$.

\subsection{Topological Equivalence of Grounding}

\begin{definition}[Grounding-Induced Topology]
The grounding-induced topology $\tau_A$ on $\Sigma$ is the initial
topology with respect to $G$: the coarsest topology making
$\sigma \mapsto G(\sigma)$ continuous with respect to the Vietoris
topology on $\mathcal{K}(S(A))$.
\end{definition}

\begin{definition}[Topological Equivalence of Grounding]
Agents $A$ and $B$ are \emph{groundingly equivalent} with respect to
$\Sigma$ if there exists a homeomorphism
$h : (\Sigma, \tau_A) \to (\Sigma, \tau_B)$.
\end{definition}

Grounding equivalence is strictly stronger than weak grounding.
Two agents may both be weakly grounded yet fail to be groundingly
equivalent. This failure is what we call a \emph{topological mismatch
of grounding}.

\subsection{Cognitive Distance}

\begin{conjecture}[Topological Sufficiency]
If agents $A$ and $B$ are groundingly equivalent with respect to
$\Sigma$, then they share every concept $\sigma \in \Sigma$ in the
full intensional sense.
\end{conjecture}

\begin{conjecture}[Topological Necessity]
If agents $A$ and $B$ share a concept $\sigma$ in the full intensional
sense, then they are groundingly equivalent with respect to $\Sigma$.
\end{conjecture}

The contrapositive of Conjecture 2 is practically useful: if the
sensory spaces of $A$ and $B$ are topologically inequivalent, then no
bijection between their symbol sets constitutes a genuine translation.

In practice, cognitive systems share concepts to varying degrees.
We therefore introduce a graded framework.

\begin{definition}[Cognitive Distance; v2 signature form]
The \emph{cognitive distance} $d_\mathrm{cog}(A, B)$ is a measure of
structural dissimilarity between the grounding-induced topologies
$\tau_A$ and $\tau_B$ over a shared symbol set~$\Sigma$.
Following the persistence signature formulation of
\cite[Def.~\ref{def:dcog-signature}]{minamoto2026topology},
$d_\mathrm{cog}$ is the tuple
\[
  d_\mathrm{cog}(A, B)
  = \Big(\,
      d_\mathrm{GH}\bigl((\Sigma, d_A), (\Sigma, d_B)\bigr),\;
      \{d_\mathrm{bn}^{(k)}(A, B)\}_{k \geq 0},\;
      \{\mathrm{PD}_k(A), \mathrm{PD}_k(B)\}_{k \geq 1}\,\Big),
\]
combining Gromov--Hausdorff distance with bottleneck distances
between $k$-th Vietoris--Rips persistence diagrams
(see Appendix~\ref{app:gh-theory} for the full construction).
In practice it is operationalized via the second-order RSA
correlation $\rho_\mathrm{cog}$ defined in
Remark~\ref{rem:rsa-operationalization}, which serves as a tractable
proxy for the $d_\mathrm{bn}^{(0)}$ component.
The three-level hierarchy of conceptual alignment is:
\end{definition}

\begin{remark}[RSA Operationalization and Second-Order Structural Similarity]
\label{rem:rsa-operationalization}
Let $\Sigma = \{\sigma_1, \ldots, \sigma_n\}$ be a fixed probe concept
set ($n = 29$ concepts spanning physical, abstract, relational,
and temporal domains; see Appendix~\ref{app:probe-set}).
For agent $A$ at hidden layer $\ell$, define the RSA dissimilarity matrix
\[
D^{A,\ell}_{ij} = 1 - \mathrm{corr}(\mathbf{h}^{A,\ell}_i,\, \mathbf{h}^{A,\ell}_j),
\]
where $\mathbf{h}^{A,\ell}_i \in \mathbb{R}^d$ is the mean hidden-state
activation at layer $\ell$ elicited by a canonical prompt for $\sigma_i$
\cite{kriegeskorte2008}.
This defines a pseudo-metric space $(\Sigma, d^{A,\ell})$.

We adopt as primary operative metric the \emph{second-order RSA correlation}
\[
\rho_\mathrm{cog}(A, B;\ell)
  = \mathrm{Spearman}\!\left(
      \mathrm{upper}(D^{A,\ell}),\,
      \mathrm{upper}(D^{B,\ell})
    \right),
\]
which measures structural similarity independently of scale.
Here $\mathrm{upper}(\cdot)$ denotes the upper-triangular entries.
$\rho_\mathrm{cog} = 1$ indicates identical relational geometry;
$\rho_\mathrm{cog} = 0$ indicates structural independence;
$\rho_\mathrm{cog} < 0$ indicates inverted relational structure.

\smallskip
\noindent\textbf{Pilot experiment (Pythia-70m).}
We apply this operationalization to three publicly available training
checkpoints of Pythia-70m \cite{biderman2023}:
step~0 (random initialization), step~1000 (early training),
and step~143000 (fully trained reference).
RSA matrices are computed at hidden layers $\ell \in \{2, 4, 6\}$
(early, middle, and final transformer layers) using $n = 29$ probe concepts.

\noindent
In all empirical analyses, $d_\mathrm{cog}$ is operationalized by
$\rho_\mathrm{cog}$; the Gromov--Hausdorff formulation serves as the
theoretical idealization whose full construction appears in
Appendix~\ref{app:gh-theory}.
More precisely, under the signature form introduced above,
$\rho_\mathrm{cog}$ is a monotone proxy for the
$d_\mathrm{bn}^{(0)}$ component of $d_\mathrm{cog}$ (merge-structure
alignment); the higher-order components $d_\mathrm{bn}^{(k)}$
($k \geq 1$) require larger concept inventories than the
twenty-nine--concept probe set to surface reliably
(see \cite[App.~\ref{app:dcog-h1}]{minamoto2026topology}), and are
outside the scope of the analyses below.

\noindent\textbf{Experiment 1: Training dynamics (Pythia-70m).}
To validate Prediction~P3, we compare three publicly available
training checkpoints of Pythia-70m \cite{biderman2023}:
step~0 (random initialization), step~1000 (early training),
and step~143000 (fully trained reference).
RSA matrices are computed at hidden layers $\ell \in \{2, 4, 6\}$
using $n = 29$ probe concepts and three prompt templates
(\textit{``The concept of \{c\} is'', ``A \{c\} is'', ``\{c\} refers to''})
averaged to reduce surface-form bias.

\begin{center}
\begin{tabular}{ccccc}
\hline
Layer $\ell$ & \multicolumn{2}{c}{$\rho_\mathrm{cog}(\cdot,\,\mathrm{ref})$}
             & \multicolumn{2}{c}{$\hat{d}_\mathrm{cog}(\cdot,\,\mathrm{ref})$} \\
             & step-0 & step-1000 & step-0 & step-1000 \\
\hline
2 & $0.095$ & $0.396$ & $0.154$ & $0.223$ \\
4 & $0.064$ & $0.317$ & $0.170$ & $0.292$ \\
6 & $0.017$ & $0.535$ & $0.176$ & $0.234$ \\
\hline
\end{tabular}
\end{center}

\noindent
$\rho_\mathrm{cog}$ increases monotonically across all layers
(step-0 $<$ step-1000, confirmed at $\ell = 2, 4, 6$),
consistent with Prediction~P3.
The near-zero variance of $D^{\mathrm{ref},\ell}$
at $\ell = 6$ ($\sigma \approx 10^{-4}$) reflects the known
representation isotropy of deeply-trained language models
\cite{ethayarajh2019}---itself a measurable signature of the
grounding-topology collapse described in Proposition~\ref{prop:overwrite}---which
renders $\hat{d}_\mathrm{cog}$ unreliable in this regime.
We therefore adopt $\rho_\mathrm{cog}$ as the primary operative metric.

\smallskip
\noindent\textbf{Experiment 2: Cross-model replication (GPT-2).}
To assess whether results are model-specific, we compute
$\rho_\mathrm{cog}$ between GPT-2 small (trained) and
Pythia-70m at both the trained and random-initialization checkpoints.

\begin{center}
\begin{tabular}{ccc}
\hline
Layer $\ell$ & GPT-2 vs Pythia-trained & GPT-2 vs Pythia-step-0 \\
\hline
2 & $0.419$ & $0.061$ \\
4 & $0.413$ & $0.128$ \\
6 & $0.419$ & $0.103$ \\
\hline
\end{tabular}
\end{center}

\noindent
Across all layers,
$\rho_\mathrm{cog}(\text{GPT-2}, \text{Pythia-trained})
  \gg \rho_\mathrm{cog}(\text{GPT-2}, \text{Pythia-step-0})$
($\approx 0.42$ vs.\ $\approx 0.10$, replicated at $\ell = 2, 4, 6$).
Two architectures trained on distinct corpora converge to similar
relational geometry over the probe concept set,
while random initialization shares no such structure.
This cross-model convergence is consistent with the claim that
training induces a shared grounding topology, rather than an
architecture-specific artifact.

\smallskip
\noindent
All code and data are available at
\url{https://github.com/minamominamoto/grounding-without-reality}.

\begin{figure}[t]
\centering
\includegraphics[width=0.92\textwidth]{cogsci_unified_fig2_rho_results.png}
\caption{
\textbf{RSA-based cognitive distance across training and models.}
\textit{(a)} $\rho_\mathrm{cog}$ between Pythia-70m checkpoints and the
fully trained reference (step-143000), measured at layers 2, 4, and 6.
All layers show monotonic increase, consistent with Prediction~P3.
\textit{(b)} Cross-model comparison: $\rho_\mathrm{cog}$ between
GPT-2 (trained) and Pythia at trained vs.\ random-initialization
checkpoints. Trained models show substantially higher structural
alignment ($\rho \approx 0.42$) than random initialization
($\rho \approx 0.10$), across all layers.
}
\label{fig:results}
\end{figure}

\smallskip
\noindent\textbf{Limitations.}
The RSA-based operationalization captures relational similarity in
activation space and serves as a proxy for grounding-induced topology.
It does not fully characterize grounding structure: results are
sensitive to prompt formulation, representation anisotropy, and
model-specific embedding geometry.
Accordingly, $\rho_\mathrm{cog}$ should be interpreted as an ordinal
measure of structural similarity rather than a complete metric of
grounding equivalence.
Extension to biological systems, where activation patterns cannot be
directly extracted, remains the primary empirical obligation of OP1.
\end{remark}

This yields a three-level hierarchy:

\begin{center}
\begin{tabular}{lll}
\hline
Level & Condition & Interpretation \\
\hline
1 & $d_\mathrm{cog} < \infty$, weak grounding & Extensional overlap only \\
2 & $d_\mathrm{cog}$ small & Approximate concept sharing \\
3 & $d_\mathrm{cog} = 0$ & Homeomorphism; full intensional sharing \\
\hline
\end{tabular}
\end{center}

% ============================================================

\section{Evidence from Comparative Cognition}
% ============================================================

The cases examined in this section are illustrative rather than
confirmatory: they motivate the need for the framework by showing that
the question of grounding structure arises in real biological systems,
but they do not by themselves establish non-homeomorphism of
grounding-induced topologies.
That determination requires the empirical program of OP1.
The cases are chosen because their sensory spaces are well characterized
and their cognitive capacities have been studied experimentally, making
them natural first test cases for the theoretical framework.

\subsection{\textit{Caenorhabditis elegans}: Minimal Grounding}

\textit{C.~elegans} has exactly 302 neurons and approximately 7,000
synapses, fully mapped \cite{white1986}.
Its sensory apparatus consists primarily of chemosensory neurons.
The organism exhibits chemotaxis---directed movement toward or away
from chemical gradients---which can be formalized as a predict-and-correct
loop: the animal generates an implicit prediction about gradient direction,
detects deviation, and adjusts locomotion \cite{bargmann2006}.
This constitutes the minimal implementation of an abductive-corrective loop.

We note a further implication: this loop does not merely implement
abduction but \emph{suspended abduction}. The organism does not commit
to a directional hypothesis on the basis of a single reading; it maintains
the hypothesis tentatively, continues to sample, and revises only when
accumulated evidence warrants it.
The 302-neuron architecture constitutes a candidate minimal biological
implementation of negative abductive capability (Section~4):
the predict-and-correct loop is structurally required to run multiple
correction cycles before the organism commits to a locomotor decision,
instantiating suspension in neural dynamics.

The sensory space $S(\textit{C.~elegans})$ is low-dimensional, dominated
by the chemosensory axis. The grounding-induced topology on whatever
proto-symbol set the organism maintains is therefore a low-dimensional
chemical topology.

\begin{example}
The symbol \textsc{food-source} in \textit{C.~elegans} is grounded in
chemical gradients. In \textit{Homo sapiens}, the same extensional
referent is grounded in a high-dimensional space including visual shape,
color, texture, weight, smell, taste, and social context. The
grounding-induced topologies are plausibly non-equivalent; the
``concept'' is therefore unlikely to be shared in the intensional sense.
\end{example}

\subsection{\textit{Drosophila melanogaster}: Intermediate Complexity}

\textit{Drosophila} possesses compound eyes sensitive to ultraviolet
and polarized light, olfactory neurons, mechanoreceptors, and gustatory
receptors.
Recent connectome data \cite{scheffer2020} reveals extensive cross-modal
convergence zones in the central complex, where olfactory, visual, and
proprioceptive information is combined nonlinearly.
This suggests $S(\textit{Drosophila})$ has a non-trivial fiber bundle
structure not captured by the product topology of Definition~1---a
concrete empirical case where the fiber bundle refinement of OP2
becomes necessary.

\subsection{\textit{Apis mellifera}: Abstract Concepts in a Non-Human Sensory Space}

The honeybee provides the most philosophically striking case.
Bees learn abstract relational concepts---``same as'' and ``different
from''---and transfer them to novel stimuli \cite{giurfa2001}.
Bees understand zero as a quantity smaller than one \cite{howard2018}
and can count up to four \cite{howard2019}.

The bees' sensory space is dominated by olfactory transduction of
volatile organic compounds, ultraviolet-sensitive photoreception,
and mechanoreception of substrate vibration.
The waggle dance encodes food-source location through bodily kinematics
calibrated to flight energy---not abstract distance
\cite{riley2005}---constituting weak grounding in the sense of
Definition~2.

The bee's concept of \textsc{same} is grounded in a sensory space
structured by olfactory, ultraviolet, and mechanoreceptive dimensions.
The human concept of \textsc{same} is grounded in a space structured
by visual, auditory, and tactile dimensions.
Both agents possess the concept in an abstract sense, but the
grounding-induced topologies plausibly differ.
On the present framework, this supports the hypothesis that the
intensional content is not shared---a conclusion that would follow if
Conjecture~2 holds.

% ============================================================
\section{Abduction as Overwrite}
% ============================================================

\subsection{Representational States and the Overwrite Transition}

Let $\mathcal{R}$ denote the space of representational states of
agent $A$.
A representational state $R \in \mathcal{R}$ encodes the agent's
current interpretation of sensory input: the active hypotheses, the
assigned categories, the weighted priors over sensory signals.

\begin{definition}[Interpretive Frame]
The interpretive frame $F(R)$ induced by state $R$ is the function
$F(R) : \mathcal{X} \to \mathcal{Y}$
that maps sensory inputs $x \in \mathcal{X}$ to interpreted outputs
$y \in \mathcal{Y}$ under the current representational state.
\end{definition}

\begin{definition}[Categorical Distinctness]
\label{def:categorical-distinctness}
Two hypotheses $h$ and $h'$ are \emph{categorially distinct} if the
partitions of $\mathcal{X}$ they induce are topologically inequivalent:
$\tau_{h'} \neq \tau_h$.
A \emph{correction} is a state change in which the induced topologies
are homeomorphic; an \emph{overwrite} is a state change in which they
are not.
\end{definition}

\begin{proposition}[Abductive Transitions are Overwrites]
\label{prop:overwrite}
Let $A$ adopt a new explanatory hypothesis $h_{t+1}$ at time $t+1$
that supersedes $h_t$.
If $h_{t+1}$ is categorially distinct from $h_t$, then the resulting
state transition is an overwrite, not a correction.
\end{proposition}

\subsection{Bias Propagation after Premature Overwrite}

\begin{definition}[Premature Overwrite]
An overwrite transition $R_t \to R_{t+1}$ is premature if the evidence
base $E_t$ is insufficient to distinguish $h_{t+1}$ from competing
hypotheses at an acceptable confidence level $\theta$:
$P(h_{t+1} \mid E_t) < \theta$ while some $h^i \neq h_{t+1}$ satisfies
$P(h^i \mid E_t) \geq \theta$.
\end{definition}

\begin{proposition}[Bias Propagation]
\label{prop:bias}
Let $R_{t+1}$ result from a premature overwrite where $h_{t+1}$ is
categorially distinct from the ground-truth hypothesis $h^*$.
Then for all subsequent inputs $x_{t'}$, $t' > t$, the interpreted
output $F(R_{t+1})(x_{t'})$ carries systematic bias toward $h_{t+1}$.
\end{proposition}

\begin{conjecture}[Non-Reducibility of Categorical Bias]
\label{conj:nonreducibility}
Within the committed frame $F(R_{t+1})$, the bias introduced by a
categorially distinct premature overwrite is non-reducible by further
Bayesian updating alone; it can only be eliminated by a subsequent
overwrite transition.
The toy model in Appendix~\ref{app:toy-model} illustrates this claim
in a two-hypothesis setting under fixed frame constraints.
Full generality across hypothesis classes and frame structures
constitutes an open problem (OP3).
\end{conjecture}

\begin{remark}[Scope]
The non-reducibility claim is restricted to categorially distinct
hypotheses. When $h_{t+1}$ and $h^*$ induce homeomorphic partitions
of $\mathcal{X}$, standard Bayesian updating will converge to the
correct posterior. The present claim applies specifically to cases in
which the adopted hypothesis restructures the topology of the
interpretive frame.
\end{remark}

\begin{corollary}
In the categorially distinct case, a premature overwrite is not
recoverable by Bayesian updating within the committed frame.
It constitutes a structural initial condition that propagates
interpretive bias until a subsequent overwrite occurs.
This corollary is conditional on the non-reducibility conjecture
in Proposition~\ref{prop:bias}; its empirical status is assessed
via Predictions P2 and P3.
\end{corollary}

\subsection{Premature Overwrite as a Dynamic Source of $d_\mathrm{cog}$ Increase}

The grounding-induced topology $\tau_A$ is not fixed; it is updated
by abductive transitions.
To make this precise, we note the following structural connection between
overwrite transitions and changes in $\tau_A$.

\begin{proposition}[Overwrite May Induce Topology Change]
\label{prop:overwrite-tau}
Let $A$ undergo an overwrite transition from state $R_t$ to $R_{t+1}$,
where the adopted hypothesis $h_{t+1}$ is categorially distinct from
$h_t$ (Definition~\ref{def:categorical-distinctness}).
Then the grounding map updates: $G_{t+1}(\sigma) \neq G_t(\sigma)$ for
at least one $\sigma \in \Sigma$.
Under the assumption that this update induces a non-homeomorphic
reorganization of symbol neighborhoods---that is, that the updated
grounding map does not happen to preserve the initial topology---the
grounding-induced topology changes: $\tau_{t+1} \neq \tau_t$,
and the two states imply positive cognitive distance:
$d_\mathrm{cog}(A_t, A_{t+1}) > 0$ under this assumption.
\end{proposition}

\begin{remark}
The conditional in Proposition~\ref{prop:overwrite-tau} is non-trivial:
it is conceivable that a categorially distinct overwrite changes the
grounding map while leaving the induced initial topology invariant, if
the new and old maps happen to generate the same open sets on $\Sigma$.
We regard this as a degenerate case unlikely to arise in systems with
rich symbol sets, but acknowledge it as a formal gap.
The assumption is satisfied in the Appendix~A toy model, where the
two hypotheses induce genuinely distinct partitions of $\mathbb{R}^2$
with non-homeomorphic induced topologies.
The general condition under which a grounding map update is guaranteed
to change $\tau_A$ is listed as OP6.
\end{remark}

Premature overwrite therefore constitutes a \emph{dynamic} source of
$d_\mathrm{cog}$ increase: the topology of the committed frame deviates
from the topology of the ideal frame, and this deviation propagates
until the next overwrite.

This connects the two parts of the paper.
Topological mismatch between agents (Section~2) has two sources:
(i) structural mismatch in sensory space, and
(ii) differential abductive history---agents that commit at different
rates or under different evidence conditions will accumulate different
topological distortions over time.
$d_\mathrm{cog}$ is therefore a dynamic quantity, not merely a static
comparison of sensory spaces.

% ============================================================
\section{Negative Abductive Capability}
% ============================================================

\subsection{Definition}

\begin{definition}[Suspension Region]
The suspension region $\mathcal{S} \subset \mathcal{R}$ is the set of
representational states in which the agent maintains active hypotheses
tentatively without committing to any single overwrite.
An agent in a suspension state accumulates evidence and updates
probability weights without triggering a frame-restructuring transition.
\end{definition}

\begin{definition}[Commitment Threshold]
The commitment threshold $\theta^* \in (0,1)$ is the minimum posterior
probability required to trigger an overwrite transition.
\end{definition}

\begin{definition}[Negative Abductive Capability (NAC)]
\label{def:nac}
An agent $A$ possesses negative abductive capability with respect to
hypothesis class $\mathcal{H}$ if:
\begin{enumerate}
  \item $A$ maintains a well-defined suspension region $\mathcal{S}$;
  \item $A$'s commitment threshold $\theta^*$ is calibrated to the
        evidence requirements of $\mathcal{H}$; and
  \item $A$ can remain in $\mathcal{S}$ for an extended observation
        window without degradation of suspension capacity.
\end{enumerate}

The term is chosen by analogy with Keats's negative capability---
the ability to remain ``in uncertainties, mysteries, doubts, without
any irritable reaching after fact and reason'' \cite{keats1817}---
recast here as a structural cognitive property rather than a poetic
temperament.
\end{definition}

\begin{remark}
NAC is not the absence of abductive capacity.
An agent with high NAC can generate hypotheses, hold them in suspension,
and commit only when evidence warrants.
An agent with low NAC commits rapidly, accumulating bias.
\end{remark}

\begin{proposition}[NAC as Structural Prerequisite]
Under the overwrite model, the expected bias accumulated per abductive
cycle is a decreasing function of NAC and is conjectured to be
independent of hypothesis-generation capacity.
This structural hypothesis implies that abductive reliability cannot
be improved by increasing generation capacity alone; NAC must be
co-present.
\end{proposition}

\medskip
\noindent
\textbf{Central thesis.}
\emph{Abduction requires negative capability.
The power to infer depends on the power to suspend inference.}

\subsection{NAC and $d_\mathrm{cog}$}

An agent with NAC calibrated to $\mathcal{H}$ remains in suspension
until $t \geq t^\dagger$ (the time at which evidence warrants
commitment).
The committed frame $F(R_{t^\dagger})$ is then homeomorphic to the
ideal frame $F(R^*)$, and $d_\mathrm{cog}(A, A^*) = 0$.
NAC is therefore a sufficient structural condition for $d_\mathrm{cog}$
minimization within the present model.

% ============================================================
\section{Empirical Predictions}
% ============================================================

The framework yields the following predictions, each connecting a
theoretical construct to a measurable quantity.

\paragraph{P1. Behavioral equivalence does not imply representational equivalence.}
Two species that demonstrate matched behavioral transfer accuracy on
an abstract relational rule (\textsc{same}, \textsc{different},
\textsc{zero}) may nevertheless differ in $d_\mathrm{cog}$.
\emph{Measurement}: Collect neural activation patterns in each species
during rule application; compute representational dissimilarity matrices
via RSA \cite{kriegeskorte2008}; compare the resulting metric spaces
via Gromov--Hausdorff distance.
\emph{Prediction}: $d_\mathrm{cog}$ will be positive and inversely
correlated with sensory space overlap across species pairs matched for
behavioral accuracy.

\paragraph{P2. NAC suppression increases premature commitment rate.}
Agents or systems whose suspension capacity is structurally suppressed
(high PCR) will commit to interpretive frames earlier under ambiguous
or gradually accumulating evidence than agents with high NAC.
\emph{Measurement (biological)}: Present ambiguous stimuli
(e.g., multistable perceptual displays; contradictory chemical gradients
for \textit{C.~elegans}) with controlled evidence accumulation rate;
measure commitment latency (time to stable behavioral choice) and
reversibility (degree to which contradictory evidence induces
re-suspension).
\emph{Measurement (AI)}: Use attribution graphs from mechanistic
interpretability tools \cite{anthropic2025circuit} to identify whether
multiple competing hypotheses are simultaneously active (suspension)
or one hypothesis suppresses others (commitment); measure the transition
latency as a function of evidence quality.
\emph{Prediction}: Systems trained with contrastive objectives will
show shorter commitment latency and lower reversibility than systems
trained with calibrated uncertainty objectives.

\paragraph{P3. Training topology determines $d_\mathrm{cog}$, not only accuracy.}
Two systems trained to the same behavioral accuracy on a grounding task
but under different training topologies (different sequences or
distributions of sensory exposures) will differ in $d_\mathrm{cog}$
relative to a reference grounding.
\emph{Measurement}: Train two models on matched accuracy targets but
with early vs.\ late presentation of informative stimuli (manipulating
premature overwrite rate); measure $d_\mathrm{cog}$ via RSA against a
fully-informed reference model.
\emph{Prediction}: Models exposed to early informative stimuli
(high premature overwrite risk) will show higher $d_\mathrm{cog}$
despite matched accuracy.

% ============================================================
\section{Implications}
% ============================================================

\subsection{Cross-Species Concept Attribution}

Behavioral experiments establish that an organism has learned a rule
(\textsc{same}, \textsc{different}, \textsc{zero}) under specific
training conditions.
The grounding map $G$ is formed through this training history.
If training conditions are informationally impoverished relative to
the full structure of the concept, the grounding-induced topology
$\tau_A$ will reflect the training experience rather than the full
intensional structure.

Two consequences follow.
First, $d_\mathrm{cog}$ measured in behavioral experiments is a lower
bound on the true cognitive distance.
Second, the grounding maps may carry the signature of premature
abduction---early hypotheses that bias all further sensory
interpretation.

Cross-species concept commensurability studies should report not only
behavioral transfer accuracy but also the \emph{training topology}:
the sequence and informativeness of sensory exposures through which
the concept was formed.

\subsection{Refining the Critique of Multimodal AI Systems}

The framework provides a more precise diagnostic for AI systems than
the standard ``grounded vs.\ ungrounded'' binary.
This is not primarily an argument about AI, but a test case for the
theory: if the framework is correct, AI systems that lack NAC should
show predictable patterns of $d_\mathrm{cog}$ elevation relative to
biological reference groundings.

Adding sensory channels to an AI system does not automatically produce
a sensory space topologically equivalent to the human sensory space.
Moreover, current multimodal systems commit to grounding assignments
via weight updates on each training batch, regardless of whether the
evidence in that batch is sufficient to determine the correct
interpretive frame.
The result is premature overwrite at scale: the model's $\tau_A$
reflects training distribution artifacts rather than sensory structure.

Current systems also lack a well-defined suspension region.
We conjecture that contrastive training objectives (e.g., CLIP-style
alignment) suppress NAC by rewarding confident cross-modal matching
rather than calibrated uncertainty; this suppression of suspension
capacity is itself a prediction of the framework (P2) and should be
tested empirically rather than assumed.

This requires distinguishing two superficially similar states that
are easy to conflate in current AI systems:

\begin{itemize}
  \item \textbf{Uncertainty}: the agent's posterior over hypotheses is
  wide---multiple hypotheses have non-negligible probability.
  This is a common state in transformer models and does not by itself
  constitute suspension.
  \item \textbf{Suspension}: the agent maintains competing hypotheses
  as active candidates \emph{without triggering frame selection}---
  the competition is reversible and is maintained across further
  evidence accumulation rather than resolved by the first confident
  signal.
  The operational sign is reversible competition maintained under
  accumulating evidence, not transient co-activation.
\end{itemize}

Uncertainty is necessary but not sufficient for suspension in the
formal sense of Definition~\ref{def:nac}.
A system that shows wide posteriors but immediately commits as soon as
any signal tips the balance has low NAC, even if its uncertainty is
calibrated.
This distinction determines what attribution graph analysis (P2) would
actually need to find: not co-activation of competing features per se,
but sustained, reversible co-activation that persists under continued
evidence accumulation.

This is a falsifiable prediction: systems with genuine suspension
capacity should show measurably higher commitment latency and reversal
rate under ambiguous stimuli than systems without it.

\subsection{Trustworthy AI: Beyond Accuracy}

A trustworthy AI system is not merely one with high accuracy.
It is one that does not prematurely fix its interpretive frame under
insufficient evidence.
We propose the following additional evaluation metrics:

\begin{enumerate}
  \item \textbf{Premature commitment rate (PCR)}: fraction of overwrite
        transitions triggered below threshold $\theta^*$.
  \item \textbf{Overwrite reversibility (OR)}: extent to which a prior
        overwrite can be undone by subsequent evidence.
  \item \textbf{Uncertainty retention (UR)}: duration over which the
        system maintains a suspension state without degradation.
  \item \textbf{Bias propagation factor (BPF)}: amplification of
        interpretive bias across successive overwrite cycles.
\end{enumerate}

% ============================================================
\section{Open Problems}
% ============================================================

The framework developed here is a theoretical proposal rather than a
completed formal theory.
The claims are advanced as structural hypotheses supported by
definitional argument and phenomenological evidence from comparative
cognition.
Their full generality across cognitive architectures, hypothesis
classes, and sensory modalities is not established and constitutes
the primary research obligation of the program.

\begin{description}
  \item[OP1.] \textbf{Measurement of grounding-induced topology.}
    How can $\tau_A$ be empirically measured for a biological or
    artificial system?
    Candidate approach: representational similarity analysis
    \cite{kriegeskorte2008} to construct a metric on $\Sigma$ from
    neural activation patterns, then compare via Gromov--Hausdorff
    distance.

  \item[OP2.] \textbf{Fiber bundle refinement.}
    Definition~1 uses the product topology, ignoring cross-modal
    coupling. A fiber bundle formulation, where the base space captures
    the dominant modality and fibers capture cross-modal modulation,
    is more biologically faithful. The \textit{Drosophila} central
    complex \cite{scheffer2020} provides a concrete test case.

  \item[OP3.] \textbf{Proof of Conjecture~2 (Topological Necessity).}
    A counterexample would require two agents with topologically
    inequivalent grounding yet identical relational structure over
    $\Sigma$.

  \item[OP4.] \textbf{Minimum neuron count for suspended abduction.}
    The predict-and-correct loop in \textit{C.~elegans} runs in 302
    neurons. What is the minimum network size to implement a suspension
    region of depth $m$ for a hypothesis class of size $k$?

  \item[OP5.] \textbf{NAC measurement.}
    For biological agents: measure commitment latency and reversibility
    under ambiguous stimuli with varying time pressure.
    For AI systems: Anthropic's circuit tracing tools
    \cite{anthropic2025circuit} may reveal whether internal states
    corresponding to maintained hypothesis competition (suspension)
    exist in trained models.

  \item[OP6.] \textbf{Formal relationship between NAC and
    $d_\mathrm{cog}$.}
    Establish a quantitative relationship between the premature
    commitment rate of a system and the resulting increase in
    $d_\mathrm{cog}$ relative to an ideally grounded reference.
\end{description}

% ============================================================
\section{Conclusion}
% ============================================================

Symbol grounding is not a binary property.
Whether a symbol system is grounded tells us only that some connection
exists between symbols and sensorimotor states---it tells us nothing
about the structure of that connection.
And the structure of that connection is not fixed: it is updated
through abductive transitions that may preserve or distort the
grounding-induced topology depending on when commitment is made.

We have introduced two formal tools for making this precise.
The cognitive distance $d_\mathrm{cog}$ measures the structural gap
between grounding-induced topologies and provides a continuous measure
of conceptual commensurability.
Negative abductive capability (NAC) characterizes the structural
capacity required to prevent premature topological distortion through
abductive overwrite.

Together, these tools yield a structural account of reliable cognition:
a system reliably shares concepts with others if its grounding-induced
topology matches theirs ($d_\mathrm{cog}$ small), and it maintains
that match over time by committing to interpretive frames only when
the evidence base warrants it (NAC sufficient, within the present
model).

The symbol grounding problem, thirty-six years after its formulation,
may need to become the \emph{symbol grounding structure problem}---
and the abduction problem may need to become the \emph{abductive
dynamics problem}.

\medskip
\noindent
\textit{Abduction requires negative capability.
The power to infer depends on the power to suspend inference.}

% ============================================================
\section*{Acknowledgments}
% ============================================================

No funding was received for this work.

% ============================================================
\appendix
\section{Formal Construction of Cognitive Distance via Gromov--Hausdorff Theory}
\label{app:gh-theory}

The formal definition of $d_\mathrm{cog}$ proceeds as follows.
Let $S(A) = \prod_{i=1}^{k} T_i$ be the sensory space of agent $A$,
equipped with the product topology from transduction channels $T_i$.
The grounding map $G: \Sigma \to \mathcal{K}(S(A))$ assigns each symbol
$\sigma \in \Sigma$ a compact subset of sensory states, endowed with the
Vietoris topology.
This induces a topology $\tau_A$ on $\Sigma$: the coarsest topology
making each $G(\sigma)$ continuous with respect to the Hausdorff metric
on $\mathcal{K}(S(A))$.

To compare two agents $A$ and $B$, one constructs compatible metrics
$d^A$ and $d^B$ on $\Sigma$ from $\tau_A$ and $\tau_B$ respectively
(OP1 in Appendix~\ref{app:gh-theory}), then defines:
\[
d_\mathrm{cog}(A, B) = d_{GH}((\Sigma, d^A), (\Sigma, d^B)),
\]
where $d_{GH}$ is the Gromov--Hausdorff distance \cite{gromov1981}.
This yields $d_\mathrm{cog} = 0$ if and only if $(\Sigma, d^A)$ and
$(\Sigma, d^B)$ are isometric---corresponding to full intensional
sharing of all concepts in $\Sigma$.

Computing exact Gromov--Hausdorff distances is NP-hard in general
\cite{memoli2007}.
The RSA-based operationalization (Remark~\ref{rem:rsa-operationalization})
provides a tractable proxy: under the assumption that the Pearson
correlation structure of hidden-state activations approximates the
relational geometry of $\tau_A$, the second-order RSA correlation
$\rho_\mathrm{cog}$ serves as an ordinal surrogate for $d_{GH}$.
The Frobenius surrogate $\hat{d}_\mathrm{cog} = \|D^A - D^B\|_F / n$
is a computable upper bound on $d_{GH}$ when activation variances are
comparable across agents, but degrades when one matrix collapses
(as observed at the final layer of deeply-trained models;
see Remark~\ref{rem:rsa-operationalization}).

\section{Probe Concept Set}
\label{app:probe-set}

The following $n = 29$ concepts constitute the fixed probe set $\Sigma$
used in the RSA operationalization (Remark~\ref{rem:rsa-operationalization}).
Concepts were selected to span four semantic domains, following the
cross-domain coverage principle of Kriegeskorte et al.\ \cite{kriegeskorte2008}.

\begin{center}
\begin{tabular}{ll}
\hline
Domain & Concepts \\
\hline
Physical / perceptual &
  cat, dog, apple, fire, water, stone, sky, tree \\
Abstract &
  justice, truth, freedom, beauty, number, cause \\
Relational &
  same, different, larger, smaller, before, after \\
Emotional / social &
  fear, trust, anger, love \\
Cognitive &
  belief, doubt, memory, knowledge, inference \\
\hline
\end{tabular}
\end{center}

\noindent
For each concept $\sigma_i$, activations are elicited using the prompt
template \texttt{``The concept of \{concept\} is''}, and the mean
last-layer hidden state over all token positions is recorded as
$\mathbf{h}^A_i$.
All implementation details and code are available at
\url{https://github.com/minamominamoto/grounding-without-reality}.

\section{Toy Model: Bayesian Overwrite and $d_\mathrm{cog}$}
\label{app:toy-model}
% ============================================================

We provide a minimal formal instantiation of Propositions~\ref{prop:overwrite}
and~\ref{prop:bias} in a two-hypothesis Bayesian setting.

\paragraph{Setup.}
Let $\mathcal{H} = \{h_0, h_1\}$.
Let $\mathcal{X} = \mathbb{R}^2$ with two regions $\mathcal{X}_0$
and $\mathcal{X}_1$, associated with $h_0$ and $h_1$ respectively.
Observations $x_1, x_2, \ldots$ are drawn from a mixture distribution.

\paragraph{Suspension and commitment.}
The agent performs Bayesian updates while
$\max_i P(h_i \mid x_1, \ldots, x_t) < \theta^*$.
At time $t^*$, the first time a hypothesis exceeds $\theta^*$, the
agent commits: the interpretive frame is set to $p(\cdot \mid h_j)$.

\paragraph{Premature overwrite.}
If $t^* < t^\dagger$ (the time at which the maximum-likelihood
hypothesis becomes stable under the full evidence sequence), the
overwrite is premature.

\paragraph{Bias propagation.}
After commitment to $h_j$ at $t^*$, the bias for subsequent input
$x_{t'}$ is:
\[
\beta(x_{t'}) = F(R_{t^*})(x_{t'}) - F(R^*)(x_{t'}).
\]
For inputs from the non-committed region $\mathcal{X}_{1-j}$,
$\beta(x_{t'}) \neq 0$ and is not correctable by Bayesian update
within the committed frame.

\paragraph{Connection to $d_\mathrm{cog}$.}
When $h_j \neq h^\dagger$ (premature commitment to the wrong
hypothesis), the committed and ideal topologies are not homeomorphic:
\[
d_\mathrm{cog}(A, A^*) = d_{GH}((\Sigma, \tau_A), (\Sigma, \tau^*)) > 0.
\]
An agent with NAC remains in suspension until $t \geq t^\dagger$,
at which point $d_\mathrm{cog}(A, A^*) = 0$.

% ============================================================
\begin{thebibliography}{99}

\bibitem{harnad1990}
Harnad, S. (1990).
The symbol grounding problem.
\textit{Physica D}, 42(1--3), 335--346.

\bibitem{peirce1903}
Peirce, C.~S. (1903).
Pragmatism as a principle and method of right thinking.
Harvard Lectures on Pragmatism.

\bibitem{harman1965}
Harman, G. (1965).
The inference to the best explanation.
\textit{Philosophical Review}, 74(1), 88--95.

\bibitem{gardenfors2000}
G\"{a}rdenfors, P. (2000).
\textit{Conceptual Spaces: The Geometry of Thought}.
MIT Press.

\bibitem{vietoris1922}
Vietoris, L. (1922).
Bereiche zweiter Ordnung.
\textit{Monatshefte f\"{u}r Mathematik und Physik}, 32, 258--280.

\bibitem{gromov1981}
Gromov, M. (1981).
\textit{Structures m\'{e}triques pour les vari\'{e}t\'{e}s
riemanniennes}.
Textes Math\'{e}matiques, 1. CEDIC/Fernand Nathan.

\bibitem{keats1817}
Keats, J. (1817).
Letter to George and Thomas Keats, 21 December 1817.
In H.~E. Rollins (Ed.),
\textit{The Letters of John Keats} (Vol.~1).
Harvard University Press, 1958.

\bibitem{white1986}
White, J.~G., Southgate, E., Thomson, J.~N., \& Brenner, S. (1986).
The structure of the nervous system of the nematode
\textit{Caenorhabditis elegans}.
\textit{Philosophical Transactions of the Royal Society B},
314(1165), 1--340.

\bibitem{bargmann2006}
Bargmann, C.~I. (2006).
Chemosensation in \textit{C.~elegans}.
\textit{WormBook}, 1--29.

\bibitem{scheffer2020}
Scheffer, L.~K., et al. (2020).
A connectome and analysis of the adult \textit{Drosophila} central
brain.
\textit{eLife}, 9, e57443.

\bibitem{giurfa2001}
Giurfa, M., Zhang, S., Jenett, A., Menzel, R., \& Srinivasan, M.~V.
(2001).
The concepts of ``sameness'' and ``difference'' in an insect.
\textit{Nature}, 410(6831), 930--933.

\bibitem{howard2018}
Howard, S.~R., Avargues-Weber, A., Garcia, J.~E., Greentree, A.~D.,
\& Dyer, A.~G. (2018).
Numerical ordering of zero in honey bees.
\textit{Science}, 360(6393), 1124--1126.

\bibitem{howard2019}
Howard, S.~R., Avargues-Weber, A., Garcia, J.~E., Greentree, A.~D.,
\& Dyer, A.~G. (2019).
Surpassing the subitizing threshold: appetitive--aversive conditioning
improves discrimination of numerosities in honeybees.
\textit{Journal of Experimental Biology}, 222(19), jeb205658.

\bibitem{riley2005}
Riley, J.~R., Greggers, U., Smith, A.~D., Reynolds, D.~R., \&
Menzel, R. (2005).
The flight paths of honeybees recruited by the waggle dance.
\textit{Nature}, 435(7039), 205--207.

\bibitem{kriegeskorte2008}
Kriegeskorte, N., Mur, M., \& Bandettini, P. (2008).
Representational similarity analysis---connecting the branches of
systems neuroscience.
\textit{Frontiers in Systems Neuroscience}, 2, 4.

\bibitem{anthropic2025circuit}
Lindsey, J., et al. (2025).
Circuit tracing: Revealing computational graphs in language models.
Anthropic Technical Report.
\url{https://transformer-circuits.pub/2025/attribution-graphs/methods.html}

\bibitem{ethayarajh2019}
Ethayarajh, K. (2019).
How contextual are contextualized word representations?
Comparing the geometry of {BERT}, {ELMo}, and {GPT-2} embeddings.
In \textit{Proceedings of the 2019 Conference on Empirical Methods
in Natural Language Processing} (EMNLP 2019), 55--65.

\bibitem{biderman2023}
Biderman, S., Schoelkopf, H., Anthony, Q., Bradley, H., Khan, O.,
Purohit, N., Muennighoff, N., Iyer, S., Zhang, J., Portes, J.,
DeSantis, T., Pyysalo, S., Gokaslan, A., Martin, J., Phang, J.,
Gitman, I., Barak, B., Gao, L., \& Leike, J. (2023).
Pythia: A suite for analyzing large language models across training and
scaling.
In \textit{Proceedings of the 40th International Conference on Machine
Learning} (ICML 2023).

\bibitem{memoli2007}
M\'emoli, F. (2007).
On the use of Gromov--Hausdorff distances for shape comparison.
In \textit{Proceedings of the Symposium on Point-Based Graphics}, 81--90.

\end{thebibliography}

\end{document}
